\section{Related Work}
\label{sec:related}

\paragraph{Herbarium digitization pipelines.}
The closest open system is \textbf{Hespi}~\cite{turnbull2025hespi}, which
chains two-level object detection (sheet components, then fields within the
primary label), a text-type classifier, a Tesseract\,+\,TrOCR ensemble,
cross-checking taxa against World Flora Online and the Australian National
Species List, and a final GPT-4o correction. On printed labels it reaches
$92.7\,\%$ text similarity, on handwritten $88.8\,\%$, and on the heterogeneous
DILLEN set $67.1\,\%$; the LLM correction alone adds $6$--$11$ points.
Wen~\etal~\cite{wen2025ocrllm} combine PaddleOCR with a text-only LLM and a
three-tier confidence system over 16 Darwin~Core fields, reaching $95.4\,\%$
field accuracy on \emph{printed plant} labels with fully local inference---an
important baseline for whether a multimodal VLM is needed at all. Both target
vascular plants and rely on botanical authorities, which are useless for fungi.
Earlier work established that OCR on a cropped label is more accurate than on
the full sheet, and that Google Cloud Vision achieves only $\sim$24\,\% exact
matches on Latin names~\cite{owen2020nlp}, motivating cross-checks with
taxonomic references.

\paragraph{OCR/HTR and document understanding.}
Modern OCR engines include PaddleOCR, Surya, docTR, and Tesseract~5;
handwriting recognition traditionally uses TrOCR and, increasingly,
general-purpose VLMs. OCR-free models such as Donut and GOT-OCR~2.0 read an
image directly to structured text. General and OCR-specialized VLMs
(Qwen3-VL, InternVL, Gemma-4, GLM-OCR, and the Datalab family
chandra/lift/surya) have rapidly become competitive on document understanding.
Our study evaluates representatives of all these families under a single
protocol.

\paragraph{Structured extraction, validation, and adaptation.}
Constrained decoding and schema-guided generation produce syntactically valid
JSON. For fungi, the authoritative references are Index Fungorum and MycoBank,
with GBIF as a cross-check~\cite{gbif}. We test a GBIF-based validator and an
LLM-based corrector, and---unlike prior single-pass pipelines---we compare
parameter-efficient fine-tuning (LoRA~\cite{hu2022lora}) against prompting and
architecture changes. To our knowledge this is the first end-to-end study on a
multilingual mycological \emph{type} corpus of this temporal breadth.

\section{Method}
\label{sec:method}

\subsection{Task and 21-field schema}
Given a single label image, the system outputs a JSON record with
\textbf{21 fields} aligned to Darwin~Core and extended for mycology, grouped as:
\emph{identifiers} (\texttt{catalog\_number});
\emph{taxonomy} (\texttt{scientific\_name},
\texttt{verbatim\_identification}, \texttt{type\_status}, \texttt{determiner},
\texttt{verbatim\_identification\_date});
\emph{collection} (\texttt{collector}, \texttt{verbatim\_collector},
\texttt{field\_number}, \texttt{verbatim\_collection\_date});
\emph{geography} (\texttt{country}, \texttt{verbatim\_locality},
\texttt{verbatim\_latitude}, \texttt{verbatim\_longitude},
\texttt{verbatim\_elevation}, \texttt{verbatim\_habitat});
\emph{host/substrate} (\texttt{verbatim\_substrate}, \texttt{organism\_remarks},
\texttt{identification\_remarks}); and
\emph{series} (\texttt{exsiccata\_title}, \texttt{exsiccata\_number}).
Both \emph{verbatim} (as-written) and \emph{normalized} variants are requested
for several fields to disentangle reading from normalization.

\subsection{Modular four-stage pipeline}
The system is modular so that individual modules can be swapped as better models
appear (Fig.~\ref{fig:pipeline}). \textbf{(A) Preprocessing / orientation:}
EXIF transpose, page-level orientation correction via a document-orientation
classifier, a conservative Grounding-DINO\,+\,docTR fallback for labels
occupying a small image fraction, and resizing to $\text{max side}=3000$.
\textbf{(B) OCR ensemble} (optional): classical engines producing tokens with
confidences. \textbf{(C) VLM normalizer:} a vision-language model consumes the
image (and optionally OCR text) and emits the 21-field JSON. \textbf{(D)
Tool-using validator} (optional): taxonomic and geographic APIs verify
normalized fields. Each stage is evaluated independently, letting us attribute
quality gains to components. Our final production configuration collapses to
stage~A followed by a domain-adapted stage~C.

\begin{figure}[t]
\centering
\footnotesize
\setlength{\fboxsep}{4pt}
\fbox{\parbox{0.92\linewidth}{\centering
\textbf{Label image (JPG)}\\[2pt]
$\downarrow$\\[1pt]
\textbf{(A) Preprocess}: EXIF, orient (PaddleX\,/\,DINO+docTR), resize\\[2pt]
$\downarrow$\\[1pt]
\textbf{(B) OCR ensemble} \emph{(optional)}: Paddle\,+\,Surya\,+\,TrOCR\\[2pt]
$\downarrow$\\[1pt]
\textbf{(C) VLM normalizer}: Qwen3-VL-8B + LoRA $\rightarrow$ JSON$_{21}$\\[2pt]
$\downarrow$\\[1pt]
\textbf{(D) Validator} \emph{(optional)}: GBIF / Index Fungorum\\[2pt]
$\downarrow$\\[1pt]
\textbf{Darwin Core record} + per-field confidence + audit log
}}
\caption{The modular pipeline. The production system uses stages A$\rightarrow$C
(domain-adapted); stages B and D are studied as ablations
(\cref{sec:experiments}).}
\label{fig:pipeline}
\end{figure}

\subsection{Ground truth, splits, and metrics}
The 394 images were manually annotated across all 21 fields (JSONL). For
fine-tuning we use a stratified split of \textbf{310 train / 79 test} images,
with 5 additional stratified images reserved as few-shot exemplars and excluded
from both. A single evaluator scores every experiment identically:
character/word error rate (CER, WER via jiwer on concatenated fields);
\emph{sim}, the RapidFuzz \texttt{token\_set\_ratio}; \emph{hard-R}, recall at a
similarity threshold $\tau=0.85$; \emph{soft-R}, mean per-field similarity where
present; $F_1$-micro (TP/FP/FN summed over all fields, our headline metric) and
$F_1$-macro; \emph{parse\_rate}, the fraction of outputs yielding valid JSON;
and \emph{sec/img}. Unless noted, ``$F_1$'' denotes $F_1$-micro.

\subsection{Fine-tuning setup}
We apply LoRA~\cite{hu2022lora} to \texttt{Qwen3-VL-8B-Instruct}. Images are
capped at $\text{max\_pixels}=1024{\times}28{\times}28$
($\approx$900$\times$900\,px) to avoid out-of-memory from native-resolution
image tokens ($\sim$6\,580 tokens otherwise). Training uses bf16, gradient
checkpointing, flash-attention~2, effective batch~8, and a cosine schedule with
$10\,\%$ warmup on a single H100~80\,GB. The final recipe, selected by the
sweep of \cref{sec:experiments}, is rank $r=32$, $\alpha=256$ ($8{\times}r$),
dropout $0.05$, $\text{lr}=7\text{e-}5$, 5 epochs, targeting \emph{all} linear
projections (attention and MLP). Training runs in $\sim$25\,min; the adapter is
334\,MB ($43.6$\,M / $8.81$\,B $=0.5\,\%$ trainable parameters).
